\documentclass[10pt]{exam}
\usepackage[phy]{template-for-exam}

\title{Egg Drop Lab}
\author{Rohrbach}
\date{\today}

\begin{document}
\maketitle

\section*{Goal:} 
  We will be dropping your egg from the balcony by the library stairwell.  Your goal is to design a device that will prevent the egg from breaking.

\section*{Materials:}

\begin{itemize}
  \item 1 egg
  \item 5 sheets of printer paper
  \item 1 meter of masking tape
\end{itemize}

\vspace{-2em}

\begin{questions}
  \uplevel{\section*{Design}}
  
    \question \label{quants}
      The egg will break if too much force is applied to it.  Therefore, your device should attempt to decrease the amount of force applied to the egg.  Consider the equation $\Sigma F = \Delta p / \Delta t$.  Which two {\bf physics quantities} affect the amount of force that the egg feels?  How does each quantity affect it?
      \vs
    
    \question
      Design a device that you think will keep the egg from breaking.  Sketch your device.
      \vs[3]
    
    \question
      Explain how your device addresses what you discussed in \#\ref{quants}.
      \vs
  
  \uplevel{\section*{Results}}
  
  
    \question
      Explain what happened.  Were you successful?
      \vspace{4em}
 \end{questions} 
  
\pagebreak
\section*{Lab Report}

On Schoology, you will turn in a Lab Report on a Word document that addresses the following:
%
\begin{enumerate}
  \item {\bf Purpose}
  \item {\bf Your design.} Explain your design and how you expected it to work.  Make sure to explain how your device intended to address momentum and/or impulse.
  \item {\bf Results.} After you drop the egg, explain what happened. Were you successful?
  \item {\bf Conclusion.}  Explain how the results differed from what you expected and what you could do differently to improve your design.  Make sure to draw a connection between your suggested improvements and the concepts of impulse or momentum.
\end{enumerate}
%
Please note, every member of your group must turn in a lab report.  This is an \emph{individual} assignment.

\section*{Grading Rubric}

\paragraph{Application to Momentum and Impulse} Your design discussion and your conclusion discussion accurately and completely draws the connection between the lab and impulse/momentum concepts.

$\circ$ 10/10 Meets Standard \hfill

$\circ$ 7/10 You address impulse/momentum, but your discussion is not quite right \hfill

$\circ$ 5/10 Impulse/Momentum is superficially addressed \hfill

$\circ$ 0/10 You do not address momentum or impulse \hfill


\paragraph{Completeness}  Your design and conclusion sections show that you have put some thought into this

$\circ$ 5/5 Meets Standard \hfill

$\circ$ 3/5 Needs work \hfill

$\circ$ 0/5 These sections are missing or extremely sparse on details \hfill



\paragraph{Bonus}  Your device was successful.

$\circ$ +3 Your egg survived \hfill
$\circ$ +0 Your egg broke \hfill{}



\paragraph{Total Score} \fillin[][5em] / 15

\end{document}